\section{Cómo leer un experimento: de la observación a la hipótesis}
El mayor valor de este tipo de clases cortas de 20 minutos casi nunca está en cubrir una fórmula más, sino en fijar un fenómeno experimental. Al repasar el material a fondo, conviene hacerse tres preguntas: qué hipótesis está poniendo a prueba el experimento, qué métrica sustituta se usó, y qué pasaría si la conclusión no se sostuviera. Solo así se evita interpretar un resultado experimental puntual como una regla general.

\begin{knowledgebox}{Tres puntos de apoyo para leer resultados experimentales}
\begin{itemize}
  \item Primero, la variable independiente: qué configuración se cambió en realidad
  \item Después, la variable observada: reward, tasa de éxito o loss
  \item Por último, el límite de extrapolación: si esta conclusión puede generalizarse a otros entornos
\end{itemize}
\end{knowledgebox}

\subsection{Resumen de la sección}
El punto central del análisis experimental no es solo «cuál fue el resultado», sino «qué juicio de entrenamiento respalda ese resultado». Este es el hábito que más vale conservar al hacer ablaciones de Agentic RL más adelante.

\newpage
